\documentclass[aspectratio=169]{beamer}
\usetheme{Frankfurt}
\usecolortheme[RGB={90,40,110}]{structure}

\title[Bayesian Inference]{A Practical Introduction to Bayesian Inference}
\subtitle{Summer methods workshop, session 3}
\author[I. Fischer]{Prof.\ Ines Fischer}
\institute[Lakeview Institute]{Lakeview Institute for Quantitative Methods}
\date{August 2026}

\begin{document}

\begin{frame}
  \titlepage
\end{frame}

\begin{frame}{Session outline}
  \tableofcontents
\end{frame}

\section{Priors and posteriors}

\begin{frame}{From prior to posterior}
  Bayes' rule updates belief in parameters $\theta$ after seeing data $y$:
  \[
    p(\theta \mid y) = \frac{p(y \mid \theta)\, p(\theta)}{p(y)}
  \]
  \begin{itemize}
    \item The prior $p(\theta)$ encodes what is plausible before the experiment.
    \item The likelihood $p(y \mid \theta)$ scores how well parameters explain the data.
    \item Weakly informative priors regularize without dominating the evidence.
  \end{itemize}
\end{frame}

\section{Computation in practice}

\begin{frame}{Making it computable}
  Posteriors rarely have closed forms, so we sample from them.
  \begin{itemize}
    \item Markov chain Monte Carlo draws correlated samples whose long-run
          distribution is the posterior.
    \item Diagnostics before trust: $\widehat{R}$ near 1, healthy effective
          sample size, and no divergent transitions.
    \item Posterior predictive checks compare simulated replicates to the
          observed data.
  \end{itemize}
\end{frame}

\begin{frame}{Before next session}
  \begin{itemize}
    \item Exercise: refit the reaction-time model with a Student-t likelihood.
    \item Reading: chapter 5 of the workshop notes, on hierarchical models.
    \item Bring one dataset from your own work for Thursday's lab.
  \end{itemize}
\end{frame}

\end{document}
